%!TEX root = ../main.tex
\chapter{Simulación}

En este capítulo se presenta la validación numérica del sistema de control
diseñado en el Capítulo 3 mediante simulaciones en MATLAB/Simulink. El objetivo
es verificar que el controlador propuesto cumple con los objetivos de control
establecidos en la \autoref{sec:objetivos_de_control}, a saber: mantener $i_d = 0$ para minimizar
pérdidas y regular $\omega_m$ hacia la referencia óptima generada por el
algoritmo MPPT, con el fin de maximizar la extracción de potencia del viento.
\\ 

Las simulaciones se organizan en dos escenarios: primero se evalúa la respuesta
del sistema ante una entrada escalón en la velocidad del viento para analizar
el comportamiento transitorio en condiciones ideales (viento constante); posteriormente, se somete
al sistema a un perfil de viento turbulento para evaluar su desempeño en
condiciones más representativas de operación real.

\section{Estructura de la simulación}

La figura \ref{fig:diag_alto_nivel} muestra la arquitectura de control de alto
nivel implementada en Simulink. El modelo se compone de tres bloques principales:
la planta (WECS), el controlador y el algoritmo MPPT. La planta implementa las
ecuaciones dinámicas del sistema acoplado \eqref{eq:full_system}, incluyendo
la dinámica mecánica del rotor y las ecuaciones eléctricas del PMSG en coordenadas
$dq$. El bloque de control implementa los controladores PI de los subsistemas 
$\Sigma_1$ y $\Sigma_2$ descritos en el lema \ref{lem:wecs}, cuya estructura en 
cascada se detalla en la figura \ref{fig:diag_lazos}.
\\

\begin{figure}[ht]
\centering
\includegraphics[width=0.85\textwidth]{block_diagram.pdf}
\caption{Arquitectura de control de alto nivel del sistema WECS. El bloque MPPT genera
la referencia de velocidad $\omega_m^*$ a partir de la potencia medida $P$ y la velocidad
del viento $v_w$. El controlador implementa los subsistemas $\Sigma_1$ y $\Sigma_2$
descritos en el Capítulo 3.}
\label{fig:diag_alto_nivel}
\end{figure}

\begin{figure}[ht]
\centering
\includegraphics[width=0.85\textwidth]{PI_strct.pdf}
\caption{Estructura de control PI con tres lazos mostrando los dos subsistemas controlados. 
El lazo externo del subsistema $\Sigma_2$ regula la velocidad angular $\omega_m$ y genera 
la referencia $i_q^\star$, el lazo interno controla $i_q$ y el lazo independiente en el 
subsistema $\Sigma_1$ controla $i_d$. Los términos de compensación no lineal corresponden 
a las expresiones \eqref{valUd} y \eqref{eqnVq}.}
\label{fig:diag_lazos}
\end{figure}

El algoritmo MPPT sigue una estrategia de perturbación y observación (P\&O),
descrita en la \autoref{sec:objetivos_de_control}. En cada paso de muestreo, el algoritmo 
compara la potencia actual $P(k)$ con la potencia del paso anterior $P(k-1)$ e incrementa
o decrementa la referencia de velocidad $\omega_m^*$ en un paso fijo $\Delta\omega$
según el signo de $\Delta P = P(k) - P(k-1)$. Para evitar que las perturbaciones
del algoritmo P\&O exciten el lazo de velocidad con variaciones abruptas, la señal
$\omega_m^*$ es filtrada mediante un filtro paso-bajas de primer orden antes de
ser enviada al controlador.

\subsection{Diseño del filtro pasabajas para la referencia de velocidad}

El algoritmo MPPT (P\&O) genera inherentemente oscilaciones en 
la señal de referencia $\omega_m^\star$ debido a sus perturbaciones periódicas. 
Para suavizar estas oscilaciones sin comprometer el desempeño del lazo de control 
de velocidad, se implementó un filtro pasabajas de primer orden:

\begin{equation}
H(s) = \frac{\omega_c}{s + \omega_c}
\end{equation}

La selección de la frecuencia de corte $\omega_c$ requiere un compromiso entre 
dos objetivos contrapuestos:

\begin{itemize}
\item Atenuar suficientemente las oscilaciones del MPPT
\item Mantener una respuesta dinámica adecuada ante cambios en la velocidad del viento
\item Preservar la separación de escalas de tiempo entre los lazos de control cascada
\end{itemize}

A partir del análisis de la respuesta del sistema en lazo cerrado, se observó 
que las oscilaciones en estado estable de $\omega_m^\star$ presentan un período 
de aproximadamente $T = 4$ s, correspondiente a una frecuencia de 
$f = 0.25$ Hz ($\omega_{osc} \approx 1.57$ rad/s).
\\

La frecuencia de corte óptima se determinó mediante un proceso iterativo:

\begin{enumerate}
\item Se estableció un rango inicial basado en la frecuencia de oscilación observada
\item Se evaluaron múltiples valores de $\omega_c$ mediante simulación
\item Se analizó el compromiso entre suavizado de la señal y preservación de la dinámica
\item Se verificó que la frecuencia seleccionada no comprometiera la estabilidad 
de los lazos internos de corriente
\end{enumerate}

Se encontró que valores de $\omega_c \geq 5$ rad/s resultan insuficientes para 
atenuar las oscilaciones (siendo significativamente mayores a $\omega_{osc}$), 
mientras que valores de $\omega_c \leq 0.2$ rad/s introducen dinámicas tan lentas 
que causan acoplamiento indeseado con los lazos de corriente, manifestándose en 
oscilaciones de $\pm 40$ A en $i_d$.
\\

El valor seleccionado de $\omega_c = 1$ rad/s, ligeramente menor que la frecuencia 
de oscilación observada, proporciona un balance adecuado entre atenuación de las 
perturbaciones del MPPT y preservación de la dinámica del sistema.

\section{Parámetros de simulación}

Las simulaciones se realizaron con los parámetros del sistema listados en la
Tabla \ref{tab:parametros_sim} y las ganancias del controlador obtenidas en la
\autoref{sec:analisis_sigmas}, resumidas en la Tabla \ref{tab:ganancias}.
\\

% ── Tabla de parámetros ──────────────────────────────────────────────────────
\begin{table}[ht]
\centering
\caption{Parámetros del sistema utilizados en la simulación.}
\label{tab:parametros_sim}
\begin{tabular}{llll}
\hline
Parámetro & Símbolo & Valor & Unidades \\ \hline
Densidad del aire        & $\rho$  & $1.225$  & kg/m$^3$ \\
Radio del rotor          & $r$     & $1.84$ & m        \\
Momento de inercia       & $J$     & $7.856$ & kg$\cdot$m$^2$ \\
Coeficiente de fricción  & $f_r$   & $0.0003035$ & N$\cdot$m$\cdot$s \\
Resistencia del estator  & $R$     & $0.3676$ & $\Omega$ \\
Inductancia              & $L$     & $3.55\times 10^{-3}$ & H        \\
Flujo magnético          & $\varphi$ & $0.2867$ & Wb     \\
Número de pares de polos & $p$     & $28$ & ---      \\ \hline
\end{tabular}
\end{table}

% ── Tabla de ganancias ───────────────────────────────────────────────────────
\begin{table}[ht]
\centering
\caption{Ganancias del controlador PI en cascada.}
\label{tab:ganancias}
\begin{tabular}{lll}
\hline
Ganancia & Valor & Lazo \\ \hline
$k_{P1}$ & $0.5$ & Corriente $i_d$ ($\Sigma_1$) \\
$k_{I1}$ & $5$ & Corriente $i_d$ ($\Sigma_1$) \\
$k_{P2}$ & $0.0355$         & Corriente $i_q$ ($\Sigma_2$) \\
$k_{I2}$ & $3.55\times10^{-4}$ & Corriente $i_q$ ($\Sigma_2$) \\
$k_{P3}$ & $50$             & Velocidad $\omega_m$ ($\Sigma_2$) \\
$k_{I3}$ & $50$             & Velocidad $\omega_m$ ($\Sigma_2$) \\ \hline
\end{tabular}
\end{table}

\section{Respuesta ante entrada escalón}

Con el objetivo de caracterizar el comportamiento transitorio del sistema en
condiciones ideales, se simuló la respuesta ante una velocidad de viento
constante de $v_w = 11$ m/s. Este escenario equivale a una entrada escalón
en el torque mecánico  y permite evaluar métricas clásicas de
desempeño del controlador de forma aislada de las perturbaciones del viento.
\\

La figura \ref{fig:omega_step} muestra la respuesta de la velocidad angular
$\omega_m$ y su referencia $\omega_m^*$ generada por el MPPT. El sistema alcanza
el estado estable en aproximadamente $t \approx 30$ s, con un
sobrepico de $\approx 6$ [rad/s]. La figura \ref{fig:corrientes_step} muestra las
corrientes $i_d$ e $i_q$: la corriente $i_d$ se regula en torno a cero durante
todo el transitorio, confirmando el correcto funcionamiento del subsistema $\Sigma_1$, mientras
que $i_q$ evoluciona para suministrar el torque eléctrico necesario para acelerar
el rotor hasta $\omega_m^*$.

\begin{figure}[ht]
    \centering
    \includegraphics[width=0.85\textwidth]{omega_step.pdf}
    \caption{Respuesta de la velocidad angular $\omega_m$ ante entrada escalón
    ($v_w = 11$ m/s constante). La línea discontinua corresponde a la referencia
    $\omega_m^*$ generada por el MPPT.}
    \label{fig:omega_step}
\end{figure}

\begin{figure}[ht]
    \centering
    \includegraphics[width=0.85\textwidth]{corrientes_step.pdf}
    \caption{Corrientes $i_d$ e $i_q$ durante la respuesta al escalón.
    La corriente $i_d$ se mantiene en torno a cero en todo momento,
    validando el objetivo de minimización de pérdidas del lazo $\Sigma_1$.}
    \label{fig:corrientes_step}
\end{figure}

\section{Respuesta ante perfil de viento turbulento}

Para evaluar la robustez del sistema ante condiciones de operación más realistas,
se simuló el sistema con el perfil de viento turbulento mostrado en la
figura \ref{fig:viento}. El perfil consiste en una fase inicial de viento constante
a $v_w = 11$ m/s durante los primeros $40$ s, seguida de variaciones estocásticas
con velocidades entre aproximadamente $9.8$ m/s y $16$ m/s.
\\

\begin{figure}[ht]
    \centering
    \includegraphics[width=0.85\textwidth]{viento.pdf}
    \caption{Perfil de velocidad del viento utilizado en la simulación. La fase
    inicial de viento constante permite al MPPT converger antes de la aplicación
    del perfil turbulento.}
    \label{fig:viento}
\end{figure}

La figura \ref{fig:omega} muestra la velocidad angular $\omega_m$ junto con la
referencia $\omega_m^*$. Durante la fase de viento constante, el sistema exhibe
una respuesta transitoria coherente con la observada en el escenario de escalón,
convergiendo a la velocidad óptima en régimen permanente. Ante las variaciones
del perfil turbulento, $\omega_m$ sigue con fidelidad las variaciones de
$\omega_m^*$, demostrando que el lazo de velocidad mantiene su desempeño ante
perturbaciones realistas.
\\

\begin{figure}[ht]
    \centering
    \includegraphics[width=0.85\textwidth]{omega.pdf}
    \caption{Velocidad angular $\omega_m$ y referencia $\omega_m^*$ ante el
    perfil de viento turbulento. Las oscilaciones en $\omega_m^*$ corresponden
    a las perturbaciones del algoritmo P\&O y son atenuadas por el filtro
    paso-bajas.}
    \label{fig:omega}
\end{figure}

Las corrientes $i_d$ e $i_q$ se muestran en la figura \ref{fig:corrientes}.
La corriente $i_d$ permanece en cero durante toda la simulación,
confirmando que $\Sigma_1$ opera correctamente incluso bajo el perfil de
viento variable. La corriente $i_q$ refleja las variaciones del torque
eléctrico requerido para seguir las fluctuaciones de $\omega_m^*$.
\\

\begin{figure}[ht]
    \centering
    \includegraphics[width=0.85\textwidth]{corrientes.pdf}
    \caption{Corrientes $i_d$ (panel superior) e $i_q$ (panel inferior) ante
    el perfil de viento turbulento.}
    \label{fig:corrientes}
\end{figure}

Finalmente, la figura \ref{fig:potencia} muestra la potencia extraída $P$
comparada con la potencia máxima teórica $P_{\max}$, calculada como:

\begin{equation}
    P_{\max}(t) = \frac{1}{2}\rho\pi r^2 v_w^3(t)\, C_p^{\max}
    \label{eq:pmax}
\end{equation}

donde $C_p^{\max} = 0.411$ es el coeficiente de potencia máximo del sistema,
obtenido en la Sección 3.1.2. Durante la fase de viento constante, $P$ converge
a $P_{\max}$ conforme el MPPT alcanza el punto de operación óptimo. En el
régimen turbulento, $P$ sigue de cerca las variaciones de $P_{\max}$, con
una brecha sistemática reducida atribuible al retardo inherente del algoritmo
P\&O y a la dinámica del lazo de velocidad. La eficiencia media de extracción
de potencia, calculada sobre el intervalo $t \geq 20$ s, es:

\begin{equation}
    \eta_{\text{avg}} = \frac{\overline{P}}{\overline{P}_{\max}}
    \times 100 = 99.8\%
    \label{eq:eficiencia}
\end{equation}

\begin{figure}[ht]
    \centering
    \includegraphics[width=0.85\textwidth]{potencia.pdf}
    \caption{Potencia extraída $P$ comparada con la potencia máxima teórica
    $P_{\max}$ ante el perfil de viento turbulento. La línea continua
    corresponde a $P_{\max}$ y la línea discontinua a $P$.}
    \label{fig:potencia}
\end{figure}
